Na ampliação de 2 de novembro de 2026, o grupo confrontou o GTFS metropolitano ARTESP, versão \texttt{20261007}, com a malha municipal detalhada IBGE de 2024, selecionando Mauá pelo código \texttt{3529401} \citep{artespgtfs2026,ibgemalha2024}. Foram incluídas as coordenadas de paradas contidas no polígono, inclusive sua borda, sem depender da presença de ``Mauá'' no nome da linha. A contagem mantém identificadores distintos, inclusive pontos em sentidos opostos; não os transforma em locais físicos únicos. Os 337 registros encontrados têm coordenadas distintas. Cinco pontos incluídos ficam a até 20 metros da divisa, e três externos também estão nessa faixa; a relação foi preservada para explicitar a sensibilidade à precisão das coordenadas e da malha.

A conferência independente dos CSV do ZIP original reproduziu os identificadores de paradas e viagens do banco importado. Contaram-se viagens distintas por \texttt{trip\_id} e data de serviço, com sequência completa de ao menos duas paradas distintas, horários válidos e não decrescentes, e embarque ou desembarque no recorte municipal. Uma viagem que passa por várias paradas foi contada uma única vez; os registros ferroviários por frequência foram contabilizados separadamente. Foram examinados 3, 7 e 8 de novembro de 2026, representando dia útil, sábado e domingo, e consultas às 08h nesses dias, além de controles na madrugada e após a vigência ARTESP. O leitor de candidatos da aplicação foi executado com os bancos reais em modo somente leitura, mantendo calendários, caminhadas e limites de busca. Os horários de consulta foram parâmetros explícitos, não observações da circulação dos veículos. A procedência, os hashes e as condições de reprodução constam do Apêndice~\ref{ap:importacao}; os resultados estão na Seção~\ref{sec:artesp-maua}.
